\subsection{Clarabel for lifted convex quadratic programs}\label{sec:method-clarabel}
Clarabel uses a homogeneous primal-dual embedding for conic programs with
a quadratic objective. The SoftQP adapter invokes the published algorithm
of Goulart and Chen (\href{https://doi.org/10.1007/s12532-026-00320-7}{primary source})
on a sparse lift. Let \(z\) contain the original decision and nonnegative
penalty auxiliaries, \(P_L\succeq0\) be its Hessian, and \(q_L\) its linear
cost. Independent lower and upper slacks represent separable penalties.
For a group-peak penalty the lift instead constrains each weighted
violation by the group epigraph variable on the appropriate side. The
penalty remains in the objective. The lifted program is
\begin{equation}\label{eq:clarabel-cone-form}
 \min_z \ \tfrac12 z^\top P_L z+q_L^\top z,\qquad
 Cz+s=b,\quad s\in\mathcal K=\{0\}^{m_e}\times\mathbb R_+^{m_i}.
\end{equation}
An equality contributes its row once to the zero-cone block. An upper
side contributes \((A_i,u_i)\); a lower side contributes
\((-A_i,-l_i)\); infinite sides contribute no rows. Here \(A_i\) is the
corresponding lifted row, and \(m_e,m_i\) count equalities and finite
inequality sides. Thus complementarity applies to inequality blocks,
while equality dual variables are free.

Write \(v\) for the cone dual and \(\tau,\kappa\) for the homogeneous
scalars. At an exact embedded solution with \(\tau>0\), the primal
and dual vectors are recovered by division by \(\tau\). Before that
recovery their stationarity and feasibility equations take the form
\begin{equation}\label{eq:clarabel-embedded-residual}
 P_Lz+C^\top v+q_L\tau=0,\qquad Cz+s-b\tau=0.
\end{equation}
The embedded Newton system also couples the homogeneous scalars to the
quadratic gap and complementarity. For the nonnegative block its
centering targets include \(s_i v_i=\mu\) and \(\tau\kappa=\mu\) for
positive barrier parameter \(\mu\). The solver computes affine and
centering directions and selects steps in the cone interiors. The
equalities are enforced in the Newton system. This describes the
embedding's role without assuming strict convexity of \(P_L\);
nonconvex curvature is outside the admitted model, and practical
termination remains a numerical observation.

\begin{algorithm}[ht]
\caption{Lift, solve and check the Clarabel candidate}\label{alg:clarabel-route}
\begin{algorithmic}[1]
\State Assemble the admitted convex lift, cone blocks and finite-side map.
\State Solve the cone form with the pinned sparse Clarabel backend.
\If{the native status is exactly \texttt{Solved}}
 \State Recover the signed interval dual and map the lifted solution to physical coordinates.
 \State Check the lifted and original-space residuals and the objective.
\Else
 \State Keep the native termination and any diagnostic direction as separate evidence.
\EndIf
\end{algorithmic}
\end{algorithm}

For the cone-side duals, accumulate \(y_i=v_{i,+}-v_{i,-}\), taking an
omitted side as zero, and copy each equality dual directly. This sign
rule gives the same \(A^\top y\) as \(C^\top v\). A direction reported
during an infeasibility termination is only a candidate for independent
certificate checks. The SDK uses Clarabel 0.11.1, upper-triangular sparse
\(P_L\), one backend thread, \(10^{-7}\) gap, feasibility and
infeasibility tolerances, and an iteration limit of \(200\,000\).
The separate tight reference route uses \(10^{-10}\) and \(10^6\)
iterations. Cost depends on finite-side assembly, sparse factor fill,
factorizations and the number of Newton steps. An exact successful
status is a prerequisite for accepting the recovered solution; lifted
auxiliary or original-space checks can still reject it.
